% p1-appB-algebraic

\section{P1 Appendix B: Algebraic Properties of the Basis Elements}

Appendix B: Algebraic Properties of the Basis Elements

                 Figure (fig-p1-appB-basis). Basis-properties triptych --- 2 / 3 / phi-pi-e columns with
                                                 independence tags.

  The key algebraic relations used throughout the paper are collected here for reference.

  Golden ratio identities. The golden ratio $\varphi$ = (1+5)/2 satisfies the minimal polynomial
  $\varphi$2 = $\varphi$ + 1, from which:

  $\varphi$-1 = $\varphi$ - 1 approx 0.61803… B.1

  $\varphi$-2 = 2 - $\varphi$ approx 0.38197… B.2

  $\varphi$-3 = 2$\varphi$ - 3 approx 0.23607… B.3

  $\varphi$2 + $\varphi$-2 = ($\varphi$ + 1) + (2 - $\varphi$) = 3. B.4

  Equation (B.4) is the Trinity identity (PhD: Ch.3 Trinity Identity). It confirms that the
  squared golden ratio and its reciprocal square sum to the integer 3, establishing the
  ternary cardinality relation central to the Lucas ring Z[$\varphi$]. In the context of this paper, it
  is used only as an algebraic completeness statement.

  Algebraic rewriting of the structural constant. Using (B.3):

  alpha$\varphi$ = ($\varphi$-3)/(2) = (2$\varphi$ - 3)/(2) = $\varphi$ - (3)/(2) = (1)/(2)($\varphi$-1)3. B.5

  This confirms that alpha$\varphi$ is a degree-3 polynomial in $\varphi$ with small integer coefficients
  --- the lowest-complexity non-trivial element of the grammar near the PDG 2024 value
  alphas(mZ) = 0.1180(9) (PDG 2024).

  Linear independence. The set {1, log 3, log $\varphi$, log pi} is treated as a working
  independence assumption over Q, not as a theorem proved here. Since log e = 1, it is
  not listed as an additional independent generator. Under this assumption, elements of
  S(C) with distinct exponent tuples map to distinct real values; if a rational relation were
  discovered, the hypothesis class would be quotient by that relation and all MDL
  penalties would be recomputed.

  Numerical values of basis elements and their logarithms:

  The log-space coordinates of the four basis elements (log 3, log $\varphi$, log pi, 1) are
  numerically pairwise incommensurable, supporting the linear-independence conjecture.

\subsection{References}

- Vasilev, D. \textit{Flos Aureus Monograph: Trinity Constants}. DOI \href{https://zenodo.org/doi/10.5281/zenodo.19227877}{10.5281/zenodo.19227877}.
- Pellis, S. "The Fine-Structure Constant and the Golden Ratio." \href{https://vixra.org/pdf/2110.0117v4.pdf}{viXra 2110.0117 v4}.
- Sherbon, M. A. "Fundamental Nature of the Fine-Structure Constant." \href{https://hal.science/hal-02341850/document}{HAL hal-02341850}.
